\newpage
\section*{Вычислительная часть}
\addcontentsline{toc}{section}{Вычислительная часть}

\subsection*{Алгоритм расчёта}
\addcontentsline{toc}{subsection}{Алгоритм расчёта}

Вычислительная часть реализует полный цикл моделирования:
\begin{enumerate}
  \item выбирается система двух массивных тел и задаётся массовый параметр $\mu$;
  \item находятся координаты всех пяти точек Лагранжа;
  \item выбирается точка $L_i$ и малое начальное возмущение;
  \item уравнения движения решаются численно;
  \item сохраняются сводка расчёта, таблица траектории, статические графики
  и интерактивная веб-страница.
\end{enumerate}

Метод Дорманда--Принса восьмого порядка -- это точный численный метод для
решения систем обыкновенных дифференциальных уравнений. Он последовательно
вычисляет состояние тела в моменты времени $t_0,t_1,t_2,\ldots$ и автоматически
контролирует локальную ошибку. В проекте этот метод используется потому, что
траектории около точек Лагранжа чувствительны к малым возмущениям, и грубый
метод интегрирования мог бы исказить вывод об устойчивости.

Для каждого запуска сохраняются не только координаты траектории, но и
производные метрики: физический масштаб в километрах и сутках, время ухода за
порог, рост отклонения, дрейф интеграла Якоби и линейная устойчивость точки.

\subsection*{Базовые сценарии}
\addcontentsline{toc}{subsection}{Базовые сценарии}

В качестве базовых примеров выбраны два запуска:
\[
  L_1:\quad x_0=x_{L_1}+10^{-4},\quad v_{x0}=v_{y0}=v_{z0}=0,
\]
\[
  L_4:\quad x_0=x_{L_4}+10^{-4},\quad v_{x0}=10^{-3},\quad v_{y0}=10^{-3}.
\]
Первый запуск показывает неустойчивость коллинеарной точки $L_1$, второй --
ограниченное движение около треугольной точки $L_4$.

\begin{table}[H]
\centering
\begin{tabular}{lcc}
\toprule
Величина & $L_1$ & $L_4$ \\
\midrule
$t_{\max}$, норм. время & 20 & 12 \\
$t_{\max}$, суток & 1162.6 & 697.6 \\
$r(0)$, норм. ед. & $1.0\cdot 10^{-4}$ & $1.0\cdot 10^{-4}$ \\
$r_{\max}$, норм. ед. & 0.018691 & 0.009679 \\
$r_{\max}$, км & $2.80\cdot 10^6$ & $1.45\cdot 10^6$ \\
Время максимума, суток & 369.6 & 697.6 \\
Порог ухода $r=0.01$ & достигнут & не достигнут \\
$t_{\rm esc}$, суток & 105.1 & -- \\
Рост $r_{\max}/r(0)$ & 186.9 & 96.8 \\
$\max|\Delta C|$ & $5.1\cdot 10^{-12}$ & $1.4\cdot 10^{-14}$ \\
Линейная устойчивость & неустойчива & устойчива \\
\bottomrule
\end{tabular}
\caption{Сводные численные характеристики базовых запусков}
\label{tab:base_metrics}
\end{table}

\subsection*{Движение около $L_1$}
\addcontentsline{toc}{subsection}{Движение около $L_1$}

На рисунке \ref{fig:l1_2d} показана траектория около $L_1$. Даже малое
начальное смещение приводит к заметному уходу от точки равновесия. На графике
$r(t)$ видно, что траектория пересекает порог $r=0.01$ примерно через
105 суток. Это соответствует теоретическому выводу: $L_1$ является
линейно неустойчивой точкой.

\begin{figure}[H]
\centering
\includegraphics[width=\textwidth]{sun_earth_l1_unstable_2d.png}
\caption{Движение около неустойчивой точки $L_1$: общий вид, локальный вид и расстояние $r(t)$}
\label{fig:l1_2d}
\end{figure}

\begin{figure}[H]
\centering
\includegraphics[width=0.82\textwidth]{sun_earth_l1_unstable_3d.png}
\caption{Трёхмерный вид траектории около $L_1$}
\label{fig:l1_3d}
\end{figure}

\subsection*{Движение около $L_4$}
\addcontentsline{toc}{subsection}{Движение около $L_4$}

На рисунке \ref{fig:l4_2d} показан устойчивый случай около $L_4$. Траектория
имеет сложную форму, однако остаётся в ограниченной окрестности треугольной
точки. За время расчёта порог ухода $r=0.01$ не достигается. Величина
$\max|\Delta C|$ также мала, что подтверждает хорошее качество численного
интегрирования.

\begin{figure}[H]
\centering
\includegraphics[width=\textwidth]{sun_earth_l4_stable_2d.png}
\caption{Движение около устойчивой точки $L_4$: общий вид, локальный вид и расстояние $r(t)$}
\label{fig:l4_2d}
\end{figure}

\begin{figure}[H]
\centering
\includegraphics[width=0.82\textwidth]{sun_earth_l4_stable_3d.png}
\caption{Трёхмерный вид траектории около $L_4$}
\label{fig:l4_3d}
\end{figure}

\subsection*{Сравнение типов движения}
\addcontentsline{toc}{subsection}{Сравнение типов движения}

Главное отличие между $L_1$ и $L_4$ состоит не только в форме траектории, но и
в механизме развития возмущения. Около $L_1$ существует неустойчивое
направление: малое смещение быстро усиливается, и тело покидает окрестность
точки. Около $L_4$ движение имеет колебательный характер: траектория может
быть вытянутой и не похожей на простую окружность, но она остаётся
ограниченной.

Для проверки этого вывода в интерактивной HTML-панели дополнительно построены:
\begin{itemize}
  \item сравнение $L_1$ и $L_4$ при одинаковом начальном смещении $dx=10^{-4}$;
  \item серия амплитуд $dx=10^{-5},10^{-4},10^{-3}$ для $L_1$ и $L_4$;
  \item сравнение шагов интегрирования $dt=0.004,0.002,0.001$;
  \item график $\Delta C(t)$ для контроля интеграла Якоби.
\end{itemize}

Эти дополнительные эксперименты показывают, что качественная разница между
коллинеарной и треугольной точкой сохраняется при изменении амплитуды
начального возмущения. Сравнение разных $dt$ показывает, что выбранный шаг
интегрирования не является источником наблюдаемого ухода от $L_1$: это свойство
самой динамической системы.
